\Needspace{12\baselineskip}
\originalchapter{Monte Carlo与有限温度平衡}
\label{chap:mc}
\originalsection{详细平衡和Metropolis采样}
统计力学期望 $\avg A=\sum_sA_s\e^{-\beta H_s}/Z$ 很难通过枚举巨大构象空间直接计算。Monte Carlo构造以该权重为稳态的随机链。若建议概率 $q(s\to s')$ 对称，接受率 $a=\min(1,\e^{-\beta\Delta H})$ 满足 $\pi(s)qa(s,s')=\pi(s')qa(s',s)$，逐对流量相等，因而目标分布稳态。非对称建议需额外乘建议概率比。

稳态正确不等于有限运行已经到达稳态。不同初始构型、加热与冷却扫描、能量和序参量的时间序列可以暴露亚稳滞留。随机链的相邻样本相关，均值方差为 $\Var(A)\,2\tau_{int}/M$，其中 $\tau_{int}=1/2+\sum_{k\ge1}\rho_A(k)$ 以采样间隔为单位。块均值误差要检查随块长是否稳定，不能把每一步都当独立样本。

\originalsection{氢吸附的格子模型}
3.320 Lab5用方格格子上的 $s_i=\pm1$ 表示占据与空位，覆盖度 $\theta=(1+\avg s)/2$。每个无序邻居对只数一次时，Hamilton量
\[
H=V_1\sum_{\langle ij\rangle}s_is_j+V_2\sum_{\langle\!\langle ij\rangle\!\rangle}s_is_j-\mu\sum_is_i,
\quad V_1=1,\ V_2=-2.
\]
最近邻排斥和对角邻居偏好同态，使半覆盖时棋盘排列能同时满足全部配对偏好。翻转一个自旋时
\[
\Delta H=-2s_i\left(V_1\sum_{j\in NN(i)}s_j+V_2\sum_{j\in NNN(i)}s_j-\mu\right).
\]
化学势项的变化是 $2\mu s_i$，不是 $-\mu s_i$。源手册某些算法示意简写会造成符号或因子混淆，这里从明确Hamilton量直接求差，并由全部能量重算核验。

无相互作用时独立单点 $\avg s=\tanh(\beta\mu)$，$\theta=[1+\tanh(\beta\mu)]/2$，所以 $\ln[\theta/(1-\theta)]=2\mu/(\kb T)$。高温与非相互作用极限要分开：如果考察固定 $\mu/T$ 而 $V_i/T\to0$，才收敛到同一Langmuir曲线。

\originalsection{有限格与热力学极限}
有限胞配分函数是一有限个指数的和，真实非解析相变只在热力学极限出现。能量峰、热容峰、覆盖度急变和棋盘序参量用于定位可能的转变，还需多尺寸比较。外场 $\mu$ 破坏占据对称性，有限扫描中的滞回不自动等于平衡两相边界。

附带代码对 $4\times4$ 周期格的全部 $2^{16}$ 个构型精确枚举，并在一组温度和化学势下与实际Metropolis采样比较；这为能量、翻转公式和覆盖度给出独立基准。此有限格结果不冒充Lab5所要求的完整无限格相图。完整Lab5章另给全部四题及实际的多尺寸相图网格、临界精化和高温比较；热力学极限和高精度共存曲线的限制仍明确保留。

\originalsection{接受率、建议与遍历}
一般Metropolis--Hastings接受率为 $a(s,s')=\min[1,\pi(s')q(s'\to s)/(\pi(s)q(s\to s'))]$，使正反边概率流相等。建议可以提高效率，但必须允许状态空间连通；一个永远不改变粒子数的链不可能采样巨正则分布。对称全子格翻转与单点翻转混合仍满足详细平衡，并可跨越某些强关联势垒。

接受率太低意味着大部分样本重复，太高也可能表示步子太小、移动缓慢。没有一个接受率保证所有物性已收敛。比较序参量的自相关、块误差和不同初态的平均，才能识别采样实际有效性。

\originalsection{涨落公式和格子模型的序参量}
从 $Z=\sum\e^{-\beta H}$ 导数得热容 $C=\operatorname{Var}(H)/(\kb T^2)$。Hamilton量含 $-\mu M$ 时，$\partial_\mu\avg M=\beta\operatorname{Var}(M)$。这些关系既是物理响应，也提供解析导数与采样统计的交叉核对。若势能本身随温度变化，还需额外导数项，不能直接套常Hamilton公式。

Lab5的棋盘序用 $m_s=N^{-1}\sum_i(-1)^{x_i+y_i}s_i$；总磁化 $m=N^{-1}\sum_i s_i$ 对应覆盖度。棋盘态 $m=0$ 却 $|m_s|=1$，只看覆盖度会漏掉有序转变。有限胞因对称切换使 $\avg m_s=0$，可用 $\avg{|m_s|}$、$\avg{m_s^2}$ 或分布；它们不是同一个热力学量。

以每对只数一次计能，空/满格每点相互作用能均为 $2V_1+2V_2=-2$，外场给 $\pm\mu$；棋盘每点能 $-2V_1+2V_2=-6$、总场贡献零。所以零温棋盘与满/空态的比较阈值为 $|\mu|=4$。有限温度边界需自由能，不能把零温能交点外推成全部相图。

\originalsection{一次失败采样怎样修正}
本卷实际最初单点Metropolis在低温棋盘/满格近共存区停留过久，与 $4\times4$ 精确枚举不一致。不是放宽误差标准让它“通过”，而是检查构型轨迹、势垒及翻转能，加入对称子格整体翻转建议，再比较独立块。由于每个子格翻转是自身逆操作且等概率建议，仍用相同能差接受率保证详细平衡。

修正后15组温度/化学势的有限格覆盖度与精确枚举比较通过预设统计容差；QA保留实际参数和误差。这个例子说明正确稳态定理无法替代有限运行的混合检查。对更大格或不同参数，仍需新的时间/尺寸收敛，不以这15个点宣称完整宏观相图完成。
